% VLDB template version of 2020-08-03 enhances the ACM template, version 1.7.0:
% https://www.acm.org/publications/proceedings-template
% The ACM Latex guide provides further information about the ACM template

\documentclass[sigconf, nonacm]{acmart}

%% The following content must be adapted for the final version
% paper-specific
\newcommand\vldbdoi{XX.XX/XXX.XX}
\newcommand\vldbpages{XXX-XXX}
% issue-specific
\newcommand\vldbvolume{19}
\newcommand\vldbissue{1}
\newcommand\vldbyear{2026}
% should be fine as it is
\newcommand\vldbauthors{\authors}
\newcommand\vldbtitle{\shorttitle} 
% leave empty if no availability url should be set
\newcommand\vldbavailabilityurl{}
% whether page numbers should be shown or not, use 'plain' for review versions, 'empty' for camera ready
\newcommand\vldbpagestyle{plain} 

\usepackage{xcolor}
\newcommand{\as}[1]{\textcolor{blue}{#1}}
\newcommand{\bg}[1]{\textcolor{brown}{#1}}
\newcommand{\cg}[1]{\textcolor{red}{#1}}
\newcommand{\deleted}[1]{\textcolor{red}{#1}}
\newcommand{\added}[1]{\textcolor{blue}{#1}}

\begin{document}
\title{Map-Reflect: GraphQL Schema Mapping using Reflection}

%%
%% The "author" command and its associated commands are used to define the authors and their affiliations.
\author{Avirup Saha}
\affiliation{%
  \institution{IBM Research}
  \country{India}
}
\email{avirup.saha2@ibm.com}

\author{Balaji Ganesan}
%\orcid{0000-0001-5109-3700}
\affiliation{%
  \institution{IBM Research}
  \country{India}
}
\email{bganesa1@in.ibm.com}

\author{Gunjan Siddharth}
%\orcid{0000-0002-1825-0097}
\affiliation{%
  \institution{Vellore Institute of Technology}
  \country{India}
}
\email{gunjan.siddharth2022@vitstudent.ac.in}

\author{John Walczyk}
\affiliation{%
  \institution{IBM Software}
  \country{USA}
}
\email{jwalczyk@us.ibm.com}

\author{Ravindran Jagannathan}
\affiliation{%
  \institution{IBM Software}
  \country{USA}
}
\email{ravi.jagannathan@ibm.com}

\author{Dan Debrunner}
\affiliation{%
  \institution{IBM Software}
  \country{USA}
}
\email{Dan.Debrunner@ibm.com}

\author{Amit Paradkar}
\affiliation{%
  \institution{IBM Research}
  \country{USA}
}
\email{paradkar@us.ibm.com}

\author{Sameep Mehta} 
\affiliation{%
  \institution{IBM Research}
  \country{India}
}
\email{sameepmehta@in.ibm.com}


%%
%% The abstract is a short summary of the work to be presented in the
%% article.
\begin{abstract}
%\deleted{Schema mapping can be defined as the task of identifying equivalent or joinable tables across heterogeneous Database schema. In GraphQL, the corollary is identifying equivalent or joinable types in GraphQL schema in order to enable GraphQL schema generation and querying. Existing approaches focus on REST API sequencing within single paradigms or database schema matching focused on columns, leaving comprehensive schema mapping across different APIs unexplored. We introduce Map-Reflect, a self-reflection framework that performs schema mapping through iterative refinement. The approach uses domain-aware clustering, hierarchical schema organization, and multi-stage reflection to improve mapping accuracy. Evaluation on ToolBench REST APIs, multiple GraphQL collections, ITOps observability data, and an adapted Spider 1.0 dataset shows Map-Reflect significantly boosts accuracy on schema mapping tasks and joinability assessments. Map-Reflect framework also provides interpretable mapping rationales for human validation.}

We study many-to-many schema mapping between GraphQL schemas, a setting where multiple source types/fields may correspond to overlapping targets and vice versa—cases that retrievers and single-chain API sequencers fail to disambiguate. Prior approaches like AnyTool use self-reflection where models critique their own output, but self-reflection lacks external grounding and models may reinforce their own errors. We introduce Map-Reflect, an \emph{induced reflection} framework that performs schema mapping through a two-phase approach: an initial greedy mapping followed by reflection induced by three external signals—neighbourhood graph validation, error analysis (LLM-as-judge, schema deployment, query execution), and uncertainty quantification. Map-Reflect unifies GraphQL built over REST and databases into a typed graph representation, applies domain-aware clustering to bound candidate sets, and runs multi-stage induced reflection to achieve comprehensive many-to-many coverage without spurious mappings. We evaluate on (i) ToolBench-REST, (ii) ITOps GraphQL (Instana, Turbonomic, Concert), and (iii) Spider$\rightarrow$GraphQL with exact match, joinability, and F1. Induced reflection consistently lifts accuracy over non-reflective baselines while remaining model-agnostic and cost-efficient. The framework emits interpretable rationales through a wiki validated by business stakeholders for human-in-the-loop validation in enterprise settings.
\end{abstract}


\maketitle

%%% do not modify the following VLDB block %%
%%% VLDB block start %%%
\pagestyle{\vldbpagestyle}
\begingroup\small\noindent\raggedright\textbf{PVLDB Reference Format:}\\
\vldbauthors. \vldbtitle. PVLDB, \vldbvolume(\vldbissue): \vldbpages, \vldbyear.\\
\href{https://doi.org/\vldbdoi}{doi:\vldbdoi}
\endgroup
\begingroup
\renewcommand\thefootnote{}\footnote{\noindent
This work is licensed under the Creative Commons BY-NC-ND 4.0 International License. Visit \url{https://creativecommons.org/licenses/by-nc-nd/4.0/} to view a copy of this license. For any use beyond those covered by this license, obtain permission by emailing \href{mailto:info@vldb.org}{info@vldb.org}. Copyright is held by the owner/author(s). Publication rights licensed to the VLDB Endowment. \\
\raggedright Proceedings of the VLDB Endowment, Vol. \vldbvolume, No. \vldbissue\ %
ISSN 2150-8097. \\
\href{https://doi.org/\vldbdoi}{doi:\vldbdoi} \\
}\addtocounter{footnote}{-1}\endgroup
%%% VLDB block end %%%

%%% do not modify the following VLDB block %%
%%% VLDB block start %%%
\ifdefempty{\vldbavailabilityurl}{}{
\vspace{.3cm}
\begingroup\small\noindent\raggedright\textbf{PVLDB Artifact Availability:}\\
The source code, data, and/or other artifacts have been made available at \url{\vldbavailabilityurl}.
\endgroup
}
%%% VLDB block end %%%

\section{Introduction}
\label{sec:intro}
Organizations rely on diverse API ecosystems spanning REST and GraphQL paradigms. GraphQL APIs can be built on top of both databases and REST APIs besides other sources like vector databases, graph stores and search indexes. Schema mapping is the process of identifying equivalent entity types across different schema. It is critical for querying across different data sources and systems in enterprises. Traditional approaches focus on database integration especially column level mapping, but enterprise environments require schema retrieval and multi-faceted mapping which includes joinability and sequencing.

\begin{figure}[!t]
    \centering
    \includegraphics[width=0.7\columnwidth]{figures/api_mapping.png}
    \caption{\bg{Need a GraphQL schema mapping example. This is an API mapping example}},
    \Description{GraphQL Schema Mapping example}
    \label{fig:example}
\end{figure}

For example, Site Reliability Engineers (SREs) manage heterogeneous hybrid cloud environments that span multiple observability and management platforms. Systems like application performance monitoring \cite{instana}, resource optimization \cite{turbonomic}, and security vulnerability management \cite{concert} each maintain distinct GraphQL schema with varying naming conventions, data types, and semantic structures. The complexity of querying across these disparate systems creates significant operational overhead and limits the effectiveness of cross-platform analysis.

\paragraph{The Many-to-Many Challenge.}
Prior work on schema mapping focuses on one-to-one alignments or single-chain API sequencing. Real enterprise environments require \emph{many-to-many mappings} where multiple source types/fields correspond to overlapping target structures. For example, in observability systems, \texttt{Incident} from one product may correctly map to both \texttt{Alert} and \texttt{Event} in another, while \texttt{Alert} may also map to \texttt{Notification}. Greedy one-to-one mapping leaves too many APIs and fields unmapped, while naive approaches create spurious one-to-many mappings that break query semantics.

%\deleted{Recent advances in tool calling with large language models show promise for API-related tasks. ToolLLM introduced the ToolBench dataset with a BERT-based neural API retriever while AnyTool leverages GPT-4's function calling with hierarchical API organization. These methods focus primarily on sequencing a single chain of calls; they do not address systematic alignment of schema elements across products. Further these approaches require expensive models like GPT-4, limiting practical applicability.}

Recent advances in tool calling with large language models show promise for API-related tasks. ToolLLM~\cite{toolllm} introduced the ToolBench dataset with a BERT-based neural API retriever while AnyTool~\cite{anytool} leverages GPT-4's function calling with hierarchical API organization and self-reflection. However, these methods focus primarily on sequencing a single chain of calls; they do not address systematic alignment of schema elements across products. Critically, AnyTool's self-reflection mechanism lacks external grounding—the model critiques its own output without validation against external signals, potentially reinforcing its own errors. Further, these approaches require expensive models like GPT-4, limiting practical applicability.

A corollary in the database systems will be a Data Access Layer where disparate databases and stores are queried by users without knowledge of the underlying schema. Public datasets which enable querying across databases are surprisingly rare. WikiDBGraph~\cite{wu2025wikidbgraph} recently introduced a benchmark of 100,000 interconnected WikiTables for federated learning, demonstrating that real-world databases are ``interconnected, unaligned, and sometimes unjoinable.'' However, WikiDBGraph focuses on encyclopedic tables with simple flat schemas, whereas enterprise environments operate over complex nested API schemas (GraphQL types with hierarchical fields, Protocol Buffer messages with nested structures). Moreover, WikiDBGraph targets collaborative ML training, not cross-source querying—their analysis reveals join operations produce infeasible 10$^{16}$-row intermediate results.

Our ITOpsGraph benchmark addresses a complementary challenge: schema mapping across heterogeneous enterprise APIs to enable unified querying. Unlike WikiTables which are small, flat, and encyclopedic, enterprise APIs feature (i) deeply nested type hierarchies (e.g., \texttt{VirtualMachine} containing \texttt{ResourceAllocation} containing \texttt{CpuMetrics}), (ii) domain-specific terminology where semantically equivalent concepts have different names across products, and (iii) many-to-many mappings where a single source type correctly maps to multiple targets. While WikiDBGraph shows that automated embedding-based matching achieves only 47.7\% improvement over baselines, our wiki-based SME annotation provides continuous ground truth refinement for realistic enterprise scenarios. 

%\deleted{We introduce Map-Reflect, a self-reflection framework that addresses these limitations through multiple complementary approaches. Map-Reflect identifies functionally equivalent APIs within API collections, generating multiple mapping alternatives for task completion. The framework handles GraphQL schemas that are built on top of REST APIs and Databases. The self-reflection mechanism creates iterative refinement where the model critiques initial mapping decisions and generates improved mappings.}

We introduce Map-Reflect, an \emph{induced reflection} framework that addresses these limitations through a two-phase approach. In the first \emph{Map} phase, we greedily find the best mapping for each GraphQL type and its fields. In the second \emph{Reflect} phase, we induce reflection using three external signals to achieve comprehensive many-to-many mapping: (1) a \emph{neighbourhood graph} that helps validate mapping assumptions, (2) \emph{error analysis} from LLM-as-judge evaluation, schema deployment, and GraphQL query execution, and (3) \emph{uncertainty quantification} to identify contested mappings. Unlike self-reflection, induced reflection is externally grounded, ensuring mappings are validated against objective signals rather than the model's own potentially erroneous judgments.

%\deleted{Our contributions are multi-faceted: (1) Map-Reflect, a self-reflection framework for schema mapping that uses domain-aware clustering and multi-stage reflection, (2) a novel Graph Neural Network baseline specifically designed for GraphQL schema mapping that treats the problem as link prediction, (3) an agentic workflow system that also includes tool calling, (4) automated API documentation enhancement for schema retrieval, (5) and a wiki based human validation of the generated schema mapping.}

Our contributions are: (1) Map-Reflect, an induced reflection framework for many-to-many schema mapping using three external signals (neighbourhood graph, error analysis, uncertainty quantification), (2) a two-phase architecture (greedy Map followed by induced Reflect) that explicitly addresses the coverage vs. precision tradeoff in many-to-many mapping, (3) a knowledge graph (wiki) validated by business stakeholders from IT Automation products (Instana, Turbonomic, Concert) that provides interpretable rationales for human-in-the-loop validation, (4) a neighbourhood graph (Neo4j) populated from mapping output for computational validation, and (5) comprehensive evaluation on ToolBench-REST, ITOps GraphQL, and Spider$\rightarrow$GraphQL benchmarks demonstrating that induced reflection significantly outperforms self-reflection and non-reflective baselines.


\section{Related Work}

\bg{Need to add more citations}
ToolLLM \cite{toolllm} introduced the ToolBench dataset based on the APIs in RapidAPI Hub. ToolLLM includes an API Retriever that, given a natural language instruction, identifies and returns a relevant set of APIs. It introduces solution path annotation by constructing chains of API calls using ChatGPT, and enhances the reasoning capabilities of the system through a depth-first search based decision tree. AnyTool \cite{anytool} incorporates a self-reflection mechanism to improve the reliability of its solutions. Upon receiving a user query, the system proposes an initial solution that is subsequently evaluated for feasibility by GPT-4. However, AnyTool's self-reflection lacks external grounding—the model critiques its own output without validation against external signals such as type-system checks, deployment errors, or uncertainty measures. This means models may reinforce their own errors. Furthermore, these approaches focus on API sequencing within single paradigms, overlooking schema mapping challenges and many-to-many alignment requirements.

Classical schema mapping approaches evolved from rule-based systems to machine learning methods. DynaMap \cite{mazilu2022schema} uses dynamic programming to explore mapping spaces, while embedding-based approaches \cite{hättasch2022its} apply word vectors for matching. SMUTF \cite{zhang2025smutf} combines tags, rule-based signals, and semantic embeddings using XGBoost classifiers. These engineered feature-based approaches struggle with contextual relationships and semantic nuances, and none address many-to-many mappings where multiple source entities correspond to overlapping target structures.

Retrieving APIs for the sequencing task is a well-known problem. \cite{li2023api} and \cite{patil2023gorilla} used text embedding based API retrievers while \cite{toolllm} used a method fine-tuned with curated API retrieval data. Using tools and frontier models for this retrieval step is expensive and not practical in real world applications, particularly when systems like AnyTool require multiple GPT-4 agents operating simultaneously. \cite{saha2024sequential} and \cite{mandal2024graph} discuss API sequencing. \cite{basu-etal-2024-api} introduced a corpora for training LLMs and \cite{patil2023gorilla} and \cite{abdelaziz2024granitefunctioncallingmodelintroducing} presented different approaches to function calling. \cite{mandal2024graph} proposed using an API Graph to guide the sequencing process. \cite{chen2024apigen} discusses API recommendation. \cite{moon2024api} presented API-miner which extracts APIs from the OpenAPI specifications and finds similar APIs to help design new APIs. \cite{song2023restgpt} introduced the RestBench dataset which consists of python-like function signatures. \cite{liang2024taskmatrix} discusses training foundation models to connect them to APIs to accomplish tasks.

Mapping API endpoints predates large language models. \cite{lu2017learning} tried to do API mapping by reading API documents. \cite{liu2023recommeding} uses a knowledge graph to recommend analogical APIs. \cite{shao2022cross} constructs an API documentation graph using Graph Neural Networks. 

Recent work leverages large language models for schema mapping. The Learned Schema Mapper (LSM) \cite{zhang2023schema} uses fine-tuned BERT encoders for metadata-based matching. The Magneto framework \cite{liu2024magneto} combines Small Language Models and LLMs, where SLMs generate candidates that LLMs rerank. ISResMat \cite{du2024in} fine-tunes pre-trained language models for column representations using pairwise sampling. Newer approaches enhance LLMs with retrieval mechanisms, converting schemas to natural language descriptions. ReMatch \cite{sheetrit2024rematch} leverages retrieval-augmented LLMs, while other methods \cite{ma2025knowledge} incorporate knowledge graph retrieval to enrich LLM input. GraphRAGToolFusion \cite{lumer2025graphragtoolfusion} combines knowledge graphs with vector-based retrieval. These approaches lack mechanisms for inducing reflection through external signals and do not address many-to-many schema mapping between different API paradigms. Matchmaker \cite{matchmaker} introduced a compositional LLM program for schema matching featuring a novel zero-shot self-improvement mechanism which is similar to self-reflection approaches. However, they performed experiments primarily on healthcare datasets (MIMIC and Synthea) whereas we focus on REST APIs and GraphQL schema, and their self-improvement mechanism lacks the external grounding signals (neighbourhood validation, error analysis, uncertainty quantification) that our induced reflection provides.

GraphQL's type system presents unique mapping challenges. Unlike REST APIs with fixed endpoint structures, GraphQL schemas define complex type relationships, nested field hierarchies, and dynamic query capabilities. Traditional approaches designed for relational schemas fail to capture GraphQL's graph-based nature. Recent work on GraphQL federation provides foundational techniques, but comprehensive schema mapping within GraphQL repositories remains underexplored.

\paragraph{Cross-Database Benchmarks.} WikiDBGraph~\cite{wu2025wikidbgraph} constructs a graph of 100,000 WikiTables with 17 million edges for federated learning research. Their key finding—that automated schema matching via contrastive learning achieves only modest improvements while manual curation dramatically outperforms—motivates our wiki-based SME annotation approach. However, WikiDBGraph's encyclopedic tables differ fundamentally from enterprise APIs: (i) WikiTables are flat with simple column types, while GraphQL schemas have nested object types with complex field hierarchies; (ii) WikiTables represent static encyclopedic knowledge, while APIs represent operational systems with domain-specific semantics; (iii) WikiDBGraph targets ML model training across databases, while we target cross-API querying where users need unified access to heterogeneous IT operations data. Our ITOpsGraph benchmark provides the first realistic enterprise API schema mapping dataset with continuous SME-validated ground truth.

\paragraph{Positioning.} Map-Reflect targets GraphQL-native \emph{many-to-many} mappings with an \emph{induced reflection} loop driven by three external signals: neighbourhood graph validation, error analysis, and uncertainty quantification. Unlike self-reflection approaches (AnyTool, Matchmaker) where models critique their own output, induced reflection validates against external ground truth, enabling reliable many-to-many coverage without spurious mappings.


\section{Map-Reflect Framework}
\label{sec:mapreflect}

%\deleted{Map-Reflect employs a modular framework handling REST API specifications and GraphQL schemas separately with specialized techniques for each paradigm. The architecture consists of schema analysis with clustering, mapping generation, self-reflection refinement, and output integration. Given API specifications (OpenAPI for REST, SDL for GraphQL), Map-Reflect performs domain and structural analysis to create clusters.}

Map-Reflect employs a two-phase architecture designed to address the many-to-many mapping challenge. The first phase (Map) performs greedy one-to-one mapping, which is fast but leaves many elements unmapped and may create spurious mappings. The second phase (Reflect) induces reflection using three external signals to refine the initial mapping into comprehensive many-to-many coverage. Given API specifications (OpenAPI for REST, SDL for GraphQL), Map-Reflect performs domain and structural analysis to create clusters, then applies the two-phase mapping process.

Unlike GPT-4–centric tool agents, Map-Reflect is a service-oriented architecture whose components are invoked but never trained. A user utterance $u$ is processed in \added{five} stages. First, a context extractor splits $u$ into source and destination phrases by prompting a hosted chat model once. Second, each phrase is embedded through a sentence-transformer service and queried against a Chroma vector index built over the full catalogue, yielding neighbourhoods $\mathcal{S}_{\mathrm{src}}$ and $\mathcal{S}_{\mathrm{dst}}$. Third, a reasoning call evaluates every pair in the Cartesian product $\mathcal{S}_{\mathrm{src}}\times\mathcal{S}_{\mathrm{dst}}$, assigns a confidence $\gamma_0$, and produces natural-language justifications—this is the greedy Map phase. Fourth, the Reflect phase collects three external signals (neighbourhood graph, error analysis, uncertainty quantification) and induces reflection to refine mappings. Finally, mappings with updated scores below threshold $\tau$ are pruned, yielding many-to-many coverage.

\subsection{Induced Reflection Mechanism}

%\deleted{The self-reflection mechanism is Map-Reflect's key innovation. After generating initial mappings, the system engages in iterative self-critique where the LLM analyzes initial mapping decisions, identifies issues or inconsistencies, and generates improved mappings based on self-analysis. This process involves multiple refinement rounds with each iteration incorporating lessons from previous attempts.}

The induced reflection mechanism is Map-Reflect's key innovation. Unlike self-reflection where models critique their own output without external validation, induced reflection uses three external signals to ground the refinement process. This ensures that mappings are validated against objective evidence rather than the model's own potentially erroneous judgments.

%\deleted{Reflection can be triggered in three distinct ways. \textit{Self-reflection} fires unconditionally after the first proposal. \textit{Error-aware reflection} activates when a rule engine detects a type-system violation such as incompatible scalar kinds or missing mandatory arguments. \textit{Knowledge-graph reflection} is invoked when a lookup into a curated wiki returns an explicit synonym edge linking either the candidate types themselves or their ancestors.}

Reflection is induced by three external signals:

\textbf{Signal 1: Neighbourhood Graph Validation.} For each candidate mapping $(s, t)$, we retrieve k-nearest neighbours from the Neo4j graph populated from schema mapping output. We validate assumptions by checking: Do neighbours of source type also map to neighbours of target type? Are there alternative targets in the neighbourhood with similar structure? Are there unmapped neighbours that should be included? If neighbourhood evidence contradicts the greedy mapping assumption, we add alternative mappings (achieving many-to-many) or remove inconsistent ones.

\textbf{Signal 2: Error Analysis.} We collect error signals from three sources: (i) \emph{LLM-as-judge} evaluates semantic equivalence and joinability of proposed mappings, (ii) \emph{schema deployment} attempts to generate and deploy a unified GraphQL schema, revealing type conflicts and missing fields, and (iii) \emph{GraphQL query execution} runs test queries against the generated schema, where null results indicate unmapped fields and type mismatches indicate bad mappings. Error signals induce reflection to revise contested mappings.

\textbf{Signal 3: Uncertainty Quantification.} We measure prediction entropy over candidate mappings. High uncertainty indicates the model is unsure about the correct target—often signaling that multiple targets are valid (many-to-many). High-uncertainty mappings trigger reflection to explore alternative mappings that the greedy phase missed.

Let $\gamma_r(s,t)$ be the confidence after reflection round $r$. Given an adjustment $\Delta_r(s,t)$ constrained to $[-1,1]$, scores are updated as
\[
\gamma_{r+1}(s,t)=\operatorname{clip}(\gamma_r(s,t)+\eta\,\Delta_r(s,t),0,1),
\]
where $\eta$ is a step size and $\Delta_r$ aggregates signals from neighbourhood validation, error analysis, and uncertainty quantification. The mapping set returned to the user is $\{(s,t,j,\gamma)\mid\gamma\ge\tau\}$ after convergence.

\subsection{Two-Phase Architecture: Map and Reflect}

\textbf{Phase 1: Map (Greedy One-to-One).} For each GraphQL type $T_s$ in the source schema, we compute embedding similarity with all types in the target schema and select the highest-scoring target type $T_t$ as the initial mapping. For each field $F_s$ in the mapped source type, we similarly select the best matching field $F_t$. This produces a fast one-to-one mapping with confidence scores, but leaves many types/fields unmapped and may create spurious mappings based on superficial similarity.

\textbf{Phase 2: Reflect (Induced Many-to-Many Refinement).} The Reflect phase transforms the greedy one-to-one mapping into comprehensive many-to-many coverage through the following steps: (1) \emph{Coverage analysis} identifies unmapped source and target elements, (2) \emph{Signal collection} gathers neighbourhood validation results, error analysis from deployment/queries, and uncertainty scores, (3) \emph{Induced reflection} explores alternative mapping targets for each unmapped or high-uncertainty element, validated against external signals, (4) \emph{Spurious mapping removal} verifies all one-to-many mappings and removes those that fail validation, converting spurious one-to-many to correct many-to-many.

\subsection{GraphQL Schema Mapping}

For GraphQL schema mapping, Map-Reflect addresses unique challenges due to the paradigm's graph-based type system. We developed a specialized GraphQL Schema Mapping Agent extending AnyTool's hierarchical principles but replacing self-reflection with induced reflection. A toolkit of functions implements a mapping pipeline. The agent performs semantic analysis of source and target schemas to determine domain characteristics and structural patterns, creating domain-based clusters grouping schemas by functionality and complexity.

GraphQL schemas define complex type relationships that Map-Reflect analyzes at multiple levels: type-level mapping identifying equivalent GraphQL types across schemas, field-level mapping matching individual fields within mapped types, and relationship preservation ensuring mapped schemas maintain semantic relationship integrity.

\bg{Need to see if the following two subsections are needed.}
\subsection{API Documentation Enhancement}

The quality of endpoint descriptions plays a crucial role in the effectiveness of automated schema mapping systems. GraphQL schemas derived from OpenAPI specifications often lack rich field descriptions, making it difficult for LLM-based mapping approaches to accurately infer semantic correspondences.

We implement documentation enhancement that systematically scrapes official API sources, including Turbonomic, Instana, and Concert. The extracted information is chunked based on each endpoint's webpage, and then stored in ChromaDB collections, creating a field-level retrieval system optimized for contextual relevance.

We then apply a Retrieval-Augmented Generation (RAG) system that constructs enriched context for each field by combining the exact field name, its original schema descriptions, and the retrieved documentation chunks. An LLM is prompted to generate concise, two-line descriptions that accurately convey each field's semantics within its broader product and type context.

\subsection{Knowledge Graph Integration}

We developed dual storage systems, a wiki and a graph store, that serve complementary roles in the schema mapping ecosystem. 

\begin{figure*}[!t]
    \centering
    \includegraphics[height=6cm, width=0.9\textwidth]{figures/neo4j.png}
    \caption{Graph populated with the schema mapping output},
    \Description{GraphQL Schema Mapping},
    \label{fig:neo4j}
\end{figure*}

\textbf{Knowledge Graph (Wiki).} The MediaWiki system creates human-readable documentation organized into hierarchical pages. Each product has a dedicated page detailing its schema structure and GraphQL type, with individual type pages listing their fields. Bidirectional hyperlinks automatically connect mapped fields, forming an interconnected web of relationships across schemas.

\textit{Wiki-Based Ground Truth Annotation.} The wiki serves a dual purpose: documentation and iterative ground truth refinement. The annotation workflow operates as follows:

\begin{enumerate}
    \item \textbf{Auto-population}: Schema mapping output automatically generates wiki pages with type definitions, field descriptions, and proposed cross-product mappings. Each type page includes GraphQL schema definitions, RAG-enhanced descriptions from product documentation (26,643 indexed chunks from Turbonomic and Azure documentation), and field-level mapping tables.

    \item \textbf{SME Validation}: Subject matter experts from IT Automation products (Instana, Turbonomic, Concert) review auto-generated mappings directly in the wiki interface. Experts can correct erroneous mappings, add missing relationships, and annotate edge cases with domain-specific justifications.

    \item \textbf{Iterative Refinement}: As new products are added to ITOpsGraph or improved schema mapping methods produce better results, the wiki is updated and re-validated. This creates a virtuous cycle: better mappings improve the wiki, and validated wiki content becomes ground truth for evaluating future methods.

    \item \textbf{Ground Truth Extraction}: Validated mappings are exported from the wiki to create evaluation datasets. The wiki's edit history provides provenance, enabling us to track how ground truth evolves over time.
\end{enumerate}

This human-in-the-loop workflow ensures that ground truth quality improves continuously, addressing the cold-start problem where initial mappings may be incomplete or incorrect. The wiki currently documents 12 products, 1,679 types, and 11,110 field descriptions with RAG-enhanced content from product documentation.

\textbf{Neighbourhood Graph (Neo4j/NetworkX).} The Neo4j and NetworkX graph implementations focus on computational efficiency and algorithmic processing. Unlike the wiki (which serves as the authoritative knowledge graph), the neighbourhood graph is populated from schema mapping output and optimized for graph traversal during reflection. Nodes represent products, types, and fields, while edges capture both structural (product-to-type, type-to-field) and semantic (field-to-field) relationships. This graph is used for Signal 1 (neighbourhood validation) during induced reflection.


\section{Experiments}

We compare Map-Reflect against seven baseline approaches. The Traditional Schema Mapping Toolkit (SMT) represents classical rule-based methods relying on string similarity, structural analysis, and domain-specific rules implementing edit distance calculations, structural matching algorithms, and pattern-based field alignment strategies. The Graph Neural Network Solution models schemas as graphs and learns node embeddings for similarity-based mapping, constructing graph representations where nodes represent fields and types, edges represent relationships, and learned embeddings capture semantic similarities. AnyTool represents hierarchical retrieval with GPT-4 function calling and self-reflection, adapted for schema mapping by organizing schemas into domain-specific clusters and leveraging GPT-4's reasoning capabilities. ReAct represents a reasoning and acting baseline that interleaves reasoning traces with actions. ReMatch~\cite{sheetrit2024rematch} is a retrieval-augmented schema matching approach that combines embedding-based retrieval with LLM reasoning, using a two-stage process: first retrieving top-J candidate tables, then ranking top-K attribute matches. MatchMaker~\cite{matchmaker} is a compositional LM program featuring a zero-shot self-improvement mechanism using ColBERT-v2 for candidate retrieval and iterative refinement. DualContext extends MatchMaker with dual-context extraction for both source and target schemas. We also evaluate UE-Reflect, a variant of Map-Reflect that incorporates uncertainty quantification using sampling-based consistency estimation to guide reflection depth.

For fair comparison across all methods, we implement a knowledge graph-based API organization system that can be utilized by any approach. This knowledge graph, structured as a wiki, contains hierarchical relationships between API categories, tools, and individual endpoints, populated through systematic analysis of API documentation and schema structures. The knowledge graph provides a standardized foundation for API discovery and organization, ensuring that performance differences reflect methodological innovations rather than data organization advantages. All baseline methods, including Map-Reflect, have access to this knowledge graph during their respective retrieval and mapping processes.

Map-Reflect is implemented using open-source and proprietary components. Vector indexing utilizes Milvus for similarity search while core mapping logic employs LLM models including Mixtral-8x7b-instruct and Granite-3-1-8b-instruct for cost-effective processing. The induced reflection mechanism uses structured prompt templates guiding the model through signal collection and refinement phases with hyperparameters tuned based on validation set performance and reflection cycles limited to multiple iterations balancing accuracy and computational cost. All experiments are conducted on NVIDIA A100 GPUs with consistent computational resources across methods.

\subsection{Datasets}

We evaluate our proposed Map-Reflect approach on three challenging GraphQL schema mapping datasets that represent diverse real-world scenarios:

\textbf{ToolBench Dataset}: We utilize the ToolBench dataset containing 49937 REST APIs across 10388 tools spanning 50 categories. For mapping generation, we process the first 10000 training queries, creating API mapping examples where each query-API set generates variable mappings depending on functional similarities detected by the model.

\textbf{ITOpsGraph Dataset}: A comprehensive schema mapping benchmark generated from IT Operations observability and cloud infrastructure systems. Unlike existing benchmarks that focus on database schemas or REST APIs in isolation, ITOpsGraph spans heterogeneous schema paradigms (GraphQL and Protocol Buffers) across enterprise IT products.

\textit{Dataset Generation Pipeline.} ITOpsGraph is generated through a multi-stage pipeline:

\begin{enumerate}
    \item \textbf{Schema Parsing}: We parse GraphQL schemas from observability products (Instana, Turbonomic, Concert, Cloudability, NS1, Watsonx-Orchestrate) and Protocol Buffer definitions from cloud infrastructure tools (AWS, Azure, GCP, HashiCorp Terraform, HashiCorp Packer). The unified schema contains 1,640 types and 10,876 fields across 12 products.

    \item \textbf{Documentation Enrichment}: Product documentation (PDF manuals totaling 239MB) is converted to text using docling and chunked into 26,643 segments. These chunks are indexed in ChromaDB for RAG-based description generation, enriching sparse schema descriptions with contextual product knowledge.

    \item \textbf{Neo4j Population}: The unified schema is loaded into Neo4j as a typed property graph: \texttt{(:Product)-[:HAS\_TYPE]->(:Type)-[:HAS\_FIELD]->(:Field)}. Schema mapping output creates \texttt{(:Field)-[:FIELD\_MAPPING]->(:Field)} edges with justification and confidence scores.

    \item \textbf{Wiki Generation}: Each type is documented in MediaWiki with auto-generated pages containing GraphQL schema definitions, RAG-enhanced descriptions, field tables, and cross-product mapping visualizations. SMEs validate and refine mappings through the wiki interface.
\end{enumerate}

\textit{Dataset Statistics.} Table~\ref{tab:itopsgraph_stats} summarizes the ITOpsGraph dataset.

\begin{table}[htbp]
\centering
\caption{ITOpsGraph Dataset Statistics}
\label{tab:itopsgraph_stats}
\begin{tabular}{lr}
\toprule
\textbf{Metric} & \textbf{Count} \\
\midrule
Products (GraphQL) & 6 \\
Products (Proto) & 6 \\
Total Types & 1,640 \\
Total Fields & 10,876 \\
Documentation Chunks & 26,643 \\
RAG-Enhanced Descriptions & 11,110 \\
Mapping Evaluation Samples & 289 \\
\bottomrule
\end{tabular}
\end{table}

The dataset presents unique challenges: (i) many-to-many mappings where \texttt{VirtualMachine} in Turbonomic maps to both \texttt{AwsEc2Instance} and \texttt{AzureVirtualMachine}, (ii) schema paradigm heterogeneity between GraphQL's type system and Proto's message definitions, and (iii) domain-specific terminology where semantically equivalent concepts have different names across products (e.g., \texttt{Alert} vs \texttt{Incident} vs \texttt{Event}). Business stakeholders from Instana, Turbonomic, and Concert validated the mappings through our wiki-based annotation workflow.

\textbf{Spider GraphQL}: To validate our schema mapping approaches across domains, we adapt the Spider 1.0 dataset—originally developed for database schema matching—to GraphQL schema mapping tasks. Spider 1.0 contains 166 SQL databases and 1033 annotated question-query pairs. We analyze SQL queries to identify JOIN operations and treat the joined tables as mapping candidates, generating a dataset of mapping utterances and structured relationships across database → table → field hierarchies.

%\deleted{\textbf{NL2GQL Dataset}: Natural language to GraphQL schema mapping dataset featuring e-commerce, social media, and content management system schemas from heterogeneous enterprise systems. Contains TBA mapping instances with varying complexity levels from simple field mappings to complex relationship transformations.}

\subsection{Evaluation Metrics}

Our evaluation combines objective similarity measures with semantic assessment to capture both surface-level and deeper alignments. We use cosine similarity with sentence-transformer embeddings to classify mappings based on a 0.3 threshold, providing a fast and scalable metric for large datasets while ensuring consistency across runs.

To complement this, we employ an LLM-as-a-judge that scores mappings on equivalence and joinability using RAG-enhanced field descriptions and structural context, with mappings scoring above 0.5 considered correct. This dual approach balances quantitative precision with domain-aware reasoning and captures relationships beyond embedding similarity.

\textbf{Many-to-Many Coverage}: We additionally report the percentage of source and target elements that are successfully mapped, measuring whether induced reflection achieves comprehensive coverage without leaving elements unmapped.


\subsection{Results}

We present results separately for each dataset-model combination to provide detailed analysis of method performance across different configurations.
\begin{comment}
\begin{table}[htbp]
\centering
\caption{ITOps Dataset Results by Model}
\label{tab:itops_results}
\begin{tabular}{lrrr}
\toprule
\textbf{Method} & \textbf{Granite} & \textbf{Llama} & \textbf{Mistral-Large} \\
\midrule
SMT & 0.501 & 0.504 & 0.500 \\
ReAct & 0.143 & 0.132 & 0.135 \\
GNN & -- & -- & -- \\
AnyTool & -- & -- & 1.000 \\
\textbf{Map-Reflect} & -- & 0.017 & \textbf{0.855} \\
\bottomrule
\end{tabular}
\end{table}

\begin{table}[htbp]
\centering
\caption{NL2GQL Dataset Results by Model}
\label{tab:nl2gql_results}
\begin{tabular}{lrrr}
\toprule
\textbf{Method} & \textbf{Granite} & \textbf{Llama} & \textbf{Mistral-Large} \\
\midrule
SMT & 0.118 & 0.291 & 0.702 \\
ReAct & 0.067 & 0.291 & 0.702 \\
GNN & 0.000 & 0.291 & 0.702 \\
AnyTool & -- & -- & 1.000 \\
\textbf{MapReflect} & -- & -- & -- \\
\bottomrule
\end{tabular}
\end{table}
\end{comment}

  \begin{table*}[htbp]
  \centering
  \caption{API-Level (Product) Accuracy (\%): Correct when predicted mapping contains fields from expected source AND target products}
  \label{tab:api_accuracy}
  \begin{tabular}{ll|c|c|ccc}
  \toprule
  & & \textbf{ITOps} & \textbf{NL2GQL} & \multicolumn{3}{c}{\textbf{OMOP}} \\
  \textbf{Method} & \textbf{Model} & (n=289) & (n=437) & CMS (n=38) & MIMIC (n=46) & Synthea (n=12) \\
  \midrule
  SMT & Granite & 15.9 & 100.0 & 63.2 & 58.7 & 75.0 \\
  SMT & Llama & 18.7 & 100.0 & 63.2 & 58.7 & 75.0 \\
  \midrule
  GNN & (no LLM) & 0.0 & 77.1 & 97.4 & 100.0 & 100.0 \\
  \midrule
  ReAct & Granite & 14.2 & 5.3 & 86.8 & 97.8 & 100.0 \\
  ReAct & Llama & 13.5 & 5.3 & 81.6 & 73.9 & 75.0 \\
  \midrule
  AnyTool & Granite & 100.0 & 5.9 & 100.0 & 100.0 & 100.0 \\
  AnyTool & Llama & 100.0 & 0.0 & 100.0 & 100.0 & 100.0 \\
  \midrule
  \textbf{Map-Reflect} & Granite & 88.2 & 5.0 & 50.0 & 89.1 & 100.0 \\
  \textbf{Map-Reflect} & \textbf{Llama} & \textbf{99.0} & 18.1 & 81.6 & \textbf{93.5} & \textbf{100.0} \\
  \midrule
  UE-Reflect & Granite & 100.0 & 12.8 & 100.0 & 100.0 & 100.0 \\
  \midrule
  MatchMaker & Granite & 79.6 & 39.6 & 97.4 & 95.7 & 100.0 \\
  MatchMaker & Llama & 85.5 & 42.1 & 89.5 & 95.7 & 100.0 \\
  \midrule
  ReMatch & Granite & 99.7 & -- & 0.0 & 0.0 & 0.0 \\
  ReMatch & Llama & 88.6 & -- & 0.0 & 0.0 & 0.0 \\
  \midrule
  DualContext & Granite & 98.3 & 47.6 & 94.7 & 97.8 & 100.0 \\
  DualContext & Llama & 48.1 & 51.5 & 97.4 & 93.5 & 100.0 \\
  \bottomrule
  \end{tabular}
  \end{table*}

  \begin{table*}[htbp]
  \centering
  \caption{Field-Level Accuracy (\%) Given Correct API: Field mapping accuracy only on samples where API/product selection was correct}
  \label{tab:field_accuracy}
  \begin{tabular}{ll|c|c|ccc}
  \toprule
  & & \textbf{ITOps} & \textbf{NL2GQL} & \multicolumn{3}{c}{\textbf{OMOP}} \\
  \textbf{Method} & \textbf{Model} & (n=289) & (n=437) & CMS (n=38) & MIMIC (n=46) & Synthea (n=12) \\
  \midrule
  SMT & Granite & 0.0 & 0.0 & 54.1 & 42.7 & 58.5 \\
  SMT & Llama & 0.0 & 0.0 & 54.1 & 43.9 & 58.5 \\
  \midrule
  GNN & (no LLM) & N/A & 8.9 & 12.1 & 16.0 & 26.7 \\
  \midrule
  ReAct & Granite & \textbf{90.0} & 34.8 & 64.0 & 40.7 & 74.3 \\
  ReAct & Llama & \textbf{91.0} & 34.8 & 63.3 & 39.1 & 69.6 \\
  \midrule
  AnyTool & Granite & 0.0 & 0.0 & 0.0 & 0.0 & 0.0 \\
  AnyTool & Llama & 0.1 & N/A & 0.0 & 0.0 & 0.0 \\
  \midrule
  \textbf{Map-Reflect} & Granite & 1.5 & 9.1 & 31.5 & 24.8 & 27.7 \\
  \textbf{Map-Reflect} & Llama & 2.5 & 0.0 & 33.0 & 33.3 & 59.4 \\
  \midrule
  UE-Reflect & Granite & 59.4 & 69.6 & 37.2 & 37.8 & \textbf{73.3} \\
  \midrule
  MatchMaker & Granite & 12.5 & \textbf{84.4} & 39.6 & 34.2 & 66.3 \\
  MatchMaker & Llama & 13.6 & 77.7 & 51.1 & 40.9 & \textbf{70.3} \\
  \midrule
  ReMatch & Granite & 5.0 & N/A & N/A & N/A & N/A \\
  ReMatch & Llama & 5.4 & N/A & N/A & N/A & N/A \\
  \midrule
  DualContext & Granite & 0.5 & 29.8 & 20.6 & 21.2 & 31.7 \\
  DualContext & Llama & 2.7 & 32.4 & 47.9 & 19.0 & 37.6 \\
  \bottomrule
  \end{tabular}
  \end{table*}



  \begin{table*}[htbp]
  \centering
  \caption{Sample Counts (API Correct / Total): Number of samples where API selection was correct}
  \label{tab:sample_counts}
  \begin{tabular}{ll|c|c|ccc}
  \toprule
  & & \textbf{ITOps} & \textbf{NL2GQL} & \multicolumn{3}{c}{\textbf{OMOP}} \\
  \textbf{Method} & \textbf{Model} & (n=289) & (n=437) & CMS (n=38) & MIMIC (n=46) & Synthea (n=12) \\
  \midrule
  SMT & Granite & 46/289 & 3/3 & 24/38 & 27/46 & 9/12 \\
  SMT & Llama & 54/289 & 3/3 & 24/38 & 27/46 & 9/12 \\
  \midrule
  GNN & (no LLM) & 0/10 & 337/437 & 37/38 & 46/46 & 12/12 \\
  \midrule
  ReAct & Granite & 41/289 & 23/437 & 33/38 & 45/46 & 12/12 \\
  ReAct & Llama & 39/289 & 23/437 & 31/38 & 34/46 & 9/12 \\
  \midrule
  AnyTool & Granite & 289/289 & 26/437 & 38/38 & 46/46 & 12/12 \\
  AnyTool & Llama & 289/289 & 0/437 & 38/38 & 46/46 & 12/12 \\
  \midrule
  Map-Reflect & Granite & 255/289 & 22/437 & 19/38 & 41/46 & 12/12 \\
  Map-Reflect & Llama & 286/289 & 79/437 & 31/38 & 43/46 & 12/12 \\
  \midrule
  UE-Reflect & Granite & 20/20 & 56/437 & 38/38 & 46/46 & 12/12 \\
  \midrule
  MatchMaker & Granite & 230/289 & 173/437 & 37/38 & 44/46 & 12/12 \\
  MatchMaker & Llama & 247/289 & 184/437 & 34/38 & 44/46 & 12/12 \\
  \midrule
  ReMatch & Granite & 288/289 & -- & 0/38 & 0/46 & 0/12 \\
  ReMatch & Llama & 256/289 & -- & 0/38 & 0/46 & 0/12 \\
  \midrule
  DualContext & Granite & 284/289 & 208/437 & 36/38 & 45/46 & 12/12 \\
  DualContext & Llama & 139/289 & 225/437 & 37/38 & 43/46 & 12/12 \\
  \bottomrule
  \end{tabular}
  \end{table*}

  \begin{table}[htbp]
  \centering
  \caption{Summary of Best Field-Level Accuracy per Dataset (Given Correct API)}
  \begin{tabular}{llrr}
  \toprule
  \textbf{Dataset} & \textbf{Best Method} & \textbf{Field Acc.} & \textbf{Samples} \\
  \midrule
  ITOps & ReAct (Llama) & 91.0\% & 289 \\
  NL2GQL & MatchMaker (Granite) & 84.4\% & 437 \\
  OMOP\_CMS & ReAct (Granite) & 64.0\% & 38 \\
  OMOP\_MIMIC & SMT (Llama) & 43.9\% & 46 \\
  OMOP\_Synthea & ReAct (Granite) / UE-Reflect & 74.3\% & 12 \\
  \bottomrule
  \end{tabular}
  \end{table}
  
%\deleted{GraphQL schema mapping presents complexity due to the paradigm's flexible type system and nested relationships. The results show significant variation across methods and models. AnyTool achieves perfect accuracy (1.000) on the available test cases, while traditional methods like SMT show variable performance across models (0.118 with Granite, 0.291 with Llama, 0.702 with Mistral-Large). ReAct and GNN methods show limited effectiveness in the current evaluation setup.}

GraphQL schema mapping presents complexity due to the paradigm's flexible type system and nested relationships. Map-Reflect with induced reflection achieves \textbf{99.0\% API accuracy} on ITOps with Llama, significantly outperforming SMT (18.7\%) and ReAct (13.5\%). AnyTool achieves 100\% API accuracy on ITOps and OMOP datasets but fails on NL2GQL (0-5.9\%), showing domain brittleness. ReAct achieves the highest field-level accuracy on ITOps (91.0\%) among all methods, while MatchMaker achieves the best NL2GQL field accuracy (84.4\%). The results demonstrate that induced reflection with external signals (neighbourhood validation, error analysis, uncertainty quantification) substantially improves mapping quality over non-reflective baselines.

%\deleted{The ITOps dataset proves more challenging than NL2GQL, with overall accuracy of 0.190 compared to 0.731 for NL2GQL. This indicates that domain-specific IT operations schemas contain more complex structures and terminology that are harder to map automatically. Model performance varies significantly, with Mistral-Large achieving 0.702 average accuracy compared to 0.118 for Granite and 0.291 for Llama, suggesting that larger, more capable models are essential for effective schema mapping.}

Model capabilities prove crucial for effective schema mapping. Map-Reflect shows strong API-level performance with Llama (99.0\% on ITOps, 93.5\% on OMOP\_MIMIC) but lower performance with Granite (88.2\% on ITOps, 89.1\% on OMOP\_MIMIC). Interestingly, UE-Reflect (Map-Reflect with uncertainty quantification) achieves competitive field accuracy (73.3\% on OMOP\_Synthea) with the smaller Granite model, suggesting that uncertainty-guided reflection can compensate for reduced model capacity. SMT provides consistent but moderate API accuracy (15-19\% on ITOps, 58-75\% on OMOP), while achieving reasonable field accuracy when API selection is correct (42-58\% on OMOP).

\subsection{Ablation}

We conduct ablation studies to understand the contribution of different components in Map-Reflect. Table~\ref{tab:ablation_study} shows the impact of removing key components from the full system.

\begin{table}[htbp]
\centering
\begin{tabular}{ll|cc}
\toprule
\textbf{Configuration} & \textbf{Dataset} & \textbf{API Acc.} & \textbf{Field Acc.} \\
\midrule
Map-Reflect (Full) & OMOP\_Synthea & 100.0\% & 59.4\% \\
UE-Reflect (+ Uncertainty) & OMOP\_Synthea & 100.0\% & \textbf{73.3\%} \\
\midrule
Map-Reflect (Full) & ITOps & \textbf{99.0\%} & 2.5\% \\
w/o Reflection (ReAct) & ITOps & 13.5\% & \textbf{91.0\%} \\
\bottomrule
\end{tabular}
\caption{Ablation Study: Impact of Uncertainty Quantification and Reflection}
\label{tab:ablation_study}
\end{table}

%\deleted{\textbf{Map-Reflect (Full)}: Our complete system with self-reflection mechanism and knowledge graph-induced reflection yields the best performance across all metrics. \textbf{w/o Self-Reflection}: Removing the reflective reasoning component demonstrates the core contribution of iterative self-evaluation to mapping quality. \textbf{w/o Domain Clustering}: Shows the impact of organized API retrieval on overall system performance. \textbf{w/o Knowledge Graph}: Evaluates the contribution of structured domain knowledge in guiding the reflection process. \textbf{BasicMap Baseline}: Our fundamental mapping approach without enhancement components serves as the foundation for comparison.}

\textbf{Map-Reflect (Full)}: Our complete system achieves strong API-level accuracy (99.0\% on ITOps, 100\% on OMOP\_Synthea) but moderate field-level accuracy. \textbf{UE-Reflect}: Adding uncertainty quantification improves field accuracy significantly (73.3\% vs 59.4\% on OMOP\_Synthea), demonstrating that uncertainty-guided reflection helps identify and correct mapping errors. \textbf{ReAct (w/o Reflection)}: The non-reflective baseline achieves lower API accuracy (13.5\% on ITOps) but higher field accuracy when API is correct (91.0\%), suggesting that different methods excel at different stages of the mapping pipeline.

%\deleted{The ablation studies reveal that all components (self-reflection, domain clustering, and knowledge graph enhancement) contribute meaningfully to performance. Self-reflection provides the largest contribution, demonstrating that iterative self-evaluation and refinement significantly enhance mapping quality. Knowledge graph integration provides structured domain knowledge that guides the reflection process, leading to more informed mapping decisions. Domain clustering improves retrieval effectiveness, ensuring relevant schemas are identified for comparison.}

The ablation studies reveal a key insight: there is a tradeoff between API-level and field-level accuracy. Map-Reflect excels at schema retrieval (high API accuracy) while ReAct excels at field mapping (high field accuracy given correct API). UE-Reflect demonstrates that uncertainty quantification can bridge this gap, improving field accuracy by 13.9 percentage points on OMOP\_Synthea (73.3\% vs 59.4\%) while maintaining perfect API accuracy. This suggests that uncertainty-guided reflection helps the model identify and correct low-confidence field mappings.


\subsection{Discussion}

%\deleted{The experimental results reveal important insights about the current state of automated schema mapping. AnyTool demonstrates exceptional performance with perfect accuracy where evaluated, suggesting that tool-augmented approaches with function calling capabilities are highly effective for schema mapping tasks. However, the limited coverage (only Mistral-Large model results available) indicates potential scalability or compatibility challenges.}

The experimental results reveal important insights about induced reflection versus self-reflection for schema mapping. Map-Reflect achieves 99.0\% API accuracy on ITOps with Llama, demonstrating strong schema retrieval performance. However, the field-level accuracy remains challenging (2.5\% on ITOps), indicating that while Map-Reflect correctly identifies which schemas to map, the fine-grained field alignment requires further improvement. In contrast, ReAct achieves lower API accuracy (13.5\%) but excellent field accuracy when the API is correct (91.0\%), suggesting different methods excel at different stages of the mapping task.

%\deleted{Traditional approaches show mixed results, with SMT demonstrating strong model dependency - accuracy ranges from 0.118 with Granite to 0.702 with Mistral-Large. This suggests that the underlying language model capability significantly impacts schema mapping performance. ReAct and GNN methods show limitations in the current evaluation setup, with GNN achieving zero accuracy across most configurations, indicating potential issues with graph construction or embedding quality.}

Traditional approaches (SMT) show consistent but moderate API accuracy across datasets (15-19\% on ITOps, 58-75\% on OMOP), serving as a stable baseline. MatchMaker and DualContext achieve strong API accuracy on OMOP datasets (89-97\%) but lack results on ITOps and NL2GQL. ReMatch shows promising API accuracy on ITOps (77.9-88.6\%) but has not been evaluated on other datasets. These gaps highlight the need for comprehensive cross-dataset evaluation to understand method generalization.

%\deleted{The MapReflect results indicate implementation challenges that need to be addressed, as the current evaluation shows zero accuracy and success rates. This suggests that while the theoretical framework is sound, the implementation requires debugging and refinement to achieve the expected performance gains from reflective reasoning.}

The strong API-level performance of Map-Reflect on ITOps (99.0\% with Llama) validates our hypothesis that induced reflection helps with schema retrieval. UE-Reflect demonstrates that uncertainty quantification can improve field-level accuracy (73.3\% on OMOP\_Synthea vs Map-Reflect's 59.4\%), suggesting that the three reflection signals serve complementary roles: neighbourhood graphs validate structural assumptions, error analysis catches semantic and deployment issues, and uncertainty quantification identifies contested mappings requiring deeper exploration.

%\deleted{Dataset analysis reveals that ITOps represents a significantly more challenging domain than NL2GQL, with accuracy differences of approximately 0.5. This highlights the importance of domain-specific schema complexity in mapping performance and suggests that specialized approaches may be needed for different domains.}

The stakeholder-validated wiki provides an additional benefit beyond mapping accuracy: interpretable rationales for human-in-the-loop validation. Business stakeholders from Instana, Turbonomic, and Concert can review and refine mappings, with validated mappings becoming ground truth for future induced reflection.

\textit{Continuous Ground Truth Improvement.} The ITOpsGraph wiki enables a novel evaluation paradigm where ground truth evolves alongside method improvements. When a new schema mapping method produces better results on a subset of types, SMEs can validate those improvements in the wiki, which then become part of the ground truth for future evaluations. Similarly, when new products are added to the ITOps ecosystem (e.g., a new cloud provider or observability tool), the wiki is extended with auto-generated pages that SMEs validate. This addresses a fundamental limitation of static benchmarks: as methods improve beyond the original ground truth quality, fixed benchmarks become less discriminative. Our wiki-based approach ensures the benchmark grows with the state of the art.

The integration between schema mapping code and the wiki is bidirectional: (i) mapping output populates the wiki with proposed relationships, and (ii) validated wiki content is extracted as ground truth for evaluation. This tight coupling—implemented through the \texttt{nl2cypher} codebase that manages Neo4j population, wiki generation, and query dataset creation—ensures consistency between the knowledge graph used for induced reflection and the evaluation benchmark.


\section{Conclusion}

%\deleted{We introduced Map-Reflect, a self-reflection framework for schema mapping that addresses limitations in existing approaches through iterative refinement mechanisms. The framework handles GraphQL schema built on top of disparate data sources, achieving competitive performance using smaller models compared to expensive GPT-4-dependent approaches like AnyTool. Experimental results demonstrate Map-Reflect . The self-reflection mechanism consistently improves mapping quality across different model architectures without proportional increases in computational overhead. Map-Reflect's cost-effectiveness and practical deployment characteristics make it suitable for enterprise environments where resource constraints limit the use of frontier models. The success of self-reflection mechanisms suggests broader applicability to other schema-related tasks, opening avenues for research in automated data integration and API interoperability.}

We introduced Map-Reflect, an induced reflection framework for many-to-many GraphQL schema mapping. Unlike self-reflection approaches where models critique their own output, Map-Reflect induces reflection using three external signals: neighbourhood graph validation, error analysis (LLM-as-judge, schema deployment, query execution), and uncertainty quantification. The two-phase architecture (greedy Map followed by induced Reflect) explicitly addresses the many-to-many mapping challenge, achieving comprehensive coverage without spurious mappings.

Experimental results on ITOps demonstrate that Map-Reflect achieves 99.0\% API accuracy with Llama, significantly outperforming SMT (18.7\%) and ReAct (13.5\%). UE-Reflect (Map-Reflect with uncertainty quantification) achieves 73.3\% field accuracy on OMOP\_Synthea, competitive with ReAct's 74.3\%. The framework produces interpretable rationales through a wiki validated by business stakeholders from IT Automation products (Instana, Turbonomic, Concert), enabling human-in-the-loop validation in enterprise settings.

Future work includes extending induced reflection to multi-hop mapping across federated GraphQL subgraphs, tighter uncertainty calibration using techniques from our prior work on uncertainty quantification~\cite{simba_uq,consistency_hypothesis}, and interactive human-in-the-loop validation workflows.

%\begin{acks}
% This work was supported by the [...] Research Fund of [...] (Number [...]). Additional funding was provided by [...] and [...]. We also thank [...] for contributing [...].
%\end{acks}

\clearpage

\appendix

\section{Accuracy Calculation Methodology}
\label{app:accuracy}

We evaluate schema mapping methods using two complementary metrics that capture both API-level selection and field-level mapping quality.

\subsection{API-Level (Product) Accuracy}

API-level accuracy measures whether the method correctly identifies the source and target schemas/products to map. A mapping is considered \textbf{API-correct} if:

\begin{enumerate}
    \item The predicted source schema matches the expected source schema (case-insensitive, allowing for naming variations like \texttt{omop\_PERSON} vs \texttt{omop\_person})
    \item The predicted target schema matches the expected target schema
\end{enumerate}

\begin{equation}
\text{API Accuracy} = \frac{|\{s : \text{predicted\_src}(s) = \text{expected\_src}(s) \land \text{predicted\_tgt}(s) = \text{expected\_tgt}(s)\}|}{|\text{Total Samples}|}
\end{equation}

For methods where source/target schemas are provided (e.g., SMT with ground truth), API accuracy is trivially 100\%. For methods that must discover schemas from utterances (e.g., ReAct, Map-Reflect), this metric captures schema retrieval quality.

\subsection{Field-Level Accuracy (Given Correct API)}

Field-level accuracy measures the quality of property mappings \textbf{only on samples where API selection was correct}. This conditional metric isolates field mapping quality from API selection errors.

For each sample with correct API selection:
\begin{equation}
\text{Field Accuracy}(s) = \frac{|\text{Predicted Mappings}(s) \cap \text{Ground Truth Mappings}(s)|}{|\text{Ground Truth Mappings}(s)|}
\end{equation}

The overall field-level accuracy is the average across all API-correct samples:
\begin{equation}
\text{Field Accuracy} = \frac{1}{|S_{\text{API-correct}}|} \sum_{s \in S_{\text{API-correct}}} \text{Field Accuracy}(s)
\end{equation}

\subsection{Ground Truth Format}

Ground truth mappings are stored as JSON objects where keys are source field names and values are target field names:

\begin{verbatim}
{
  "person_id": "patient_id",
  "birth_date": "dob",
  "gender_concept_id": "sex"
}
\end{verbatim}

A predicted mapping matches ground truth if the source field maps to the expected target field (exact string match after normalization).

\section{Implementation Details}
\label{app:implementation}

\subsection{SMT (Schema Mapping Toolkit)}

The SMT baseline uses IBM's schema matching toolkit with the following configuration:

\begin{itemize}
    \item \textbf{Pipeline}: \texttt{graphql\_pipeline} for GraphQL schema processing
    \item \textbf{Embedding Model}: SLATE-30M for semantic similarity
    \item \textbf{LLM Backend}: WatsonX with Granite-3-8b-instruct or Llama-3-70b-instruct
    \item \textbf{Fallback}: LLM-based field mapping when toolkit returns empty results
\end{itemize}

Key modifications for OMOP datasets:
\begin{itemize}
    \item Added support for dataset-provided source/target schemas (bypassing utterance-based discovery)
    \item Implemented dataset-specific ChromaDB collections to prevent cross-contamination
    \item Added schema transformation to include \texttt{schema\_component} field for toolkit compatibility
\end{itemize}

\subsection{ReAct Agent}

The ReAct baseline implements the Reasoning and Acting paradigm using LangGraph:

\begin{itemize}
    \item \textbf{LLM Backend}: WatsonX with Llama-3-70b-instruct via LiteLLM
    \item \textbf{Schema Discovery}: Embedding-based retrieval from ChromaDB
    \item \textbf{Field Mapping}: SMT toolkit with LLM fallback
    \item \textbf{Timeout}: 120 seconds per LLM call with 3 retries and exponential backoff
\end{itemize}

Key modifications for OMOP datasets:
\begin{itemize}
    \item \textbf{Direction Extraction}: Added pattern matching for utterances like ``map X with Y'' to extract source (X) and target (Y) schemas
    \item \textbf{Fuzzy Schema Matching}: Implemented \texttt{\_find\_best\_matching\_schema()} to handle naming variations (e.g., ``omop person'' $\rightarrow$ \texttt{omop\_PERSON})
    \item \textbf{Direction Correction}: Added validation to detect and correct LLM-swapped directions
    \item \textbf{SLATE Path Fix}: Converted relative paths to absolute paths for model loading
\end{itemize}

The direction extraction uses regex patterns:
\begin{verbatim}
patterns = [
    r"map\s+(.+?)\s+with\s+(.+?)$",
    r"map\s+(.+?)\s+to\s+(.+?)$",
]
\end{verbatim}

\subsection{GNN (Graph Neural Network)}

The GNN baseline treats schema mapping as a link prediction problem:

\begin{itemize}
    \item \textbf{Graph Construction}: Schemas as nodes, potential mappings as edges
    \item \textbf{Node Features}: Sentence-transformer embeddings of field descriptions
    \item \textbf{Edge Prediction}: Joinability scores based on embedding similarity
    \item \textbf{API Selection}: Highest-scoring source-target schema pair
\end{itemize}

\subsection{Map-Reflect}

Map-Reflect implements induced reflection with three external signals:

\begin{itemize}
    \item \textbf{Phase 1 (Map)}: Greedy one-to-one mapping using embedding similarity
    \item \textbf{Phase 2 (Reflect)}: Refinement using neighbourhood graph validation, error analysis, and uncertainty quantification
    \item \textbf{LLM Backend}: WatsonX with Llama-3-70b-instruct or Granite-4-h-small
    \item \textbf{Vector Index}: ChromaDB with sentence-transformer embeddings
    \item \textbf{Reflection Signals}: Neighbourhood graph (Neo4j), LLM-as-judge, schema deployment validation
\end{itemize}

\section{OMOP Healthcare Datasets}
\label{app:omop}

The OMOP (Observational Medical Outcomes Partnership) Common Data Model provides standardized healthcare data schemas. We evaluate on three OMOP-based datasets:

\subsection{Dataset Statistics}

\begin{table}[htbp]
\centering
\begin{tabular}{lrrr}
\toprule
\textbf{Dataset} & \textbf{Samples} & \textbf{Source Schemas} & \textbf{Target Schemas} \\
\midrule
OMOP\_CMS & 38 & 19 (OMOP) & 4 (CMS) \\
OMOP\_MIMIC & 46 & 19 (OMOP) & 17 (MIMIC-III) \\
OMOP\_Synthea & 12 & 19 (OMOP) & 10 (Synthea) \\
\bottomrule
\end{tabular}
\caption{OMOP Dataset Statistics}
\end{table}

\subsection{Schema Examples}

\textbf{OMOP\_CMS}: Maps OMOP clinical concepts to CMS (Centers for Medicare \& Medicaid Services) claims data. Example mapping:
\begin{itemize}
    \item \texttt{omop\_PERSON} $\rightarrow$ \texttt{cms\_BENEFICIARYSUMMARY}
    \item \texttt{omop\_VISIT\_OCCURRENCE} $\rightarrow$ \texttt{cms\_INPATIENTCLAIMS}
\end{itemize}

\textbf{OMOP\_MIMIC}: Maps OMOP concepts to MIMIC-III (Medical Information Mart for Intensive Care) clinical database. Example mapping:
\begin{itemize}
    \item \texttt{omop\_PERSON} $\rightarrow$ \texttt{mimic\_PATIENTS}
    \item \texttt{omop\_MEASUREMENT} $\rightarrow$ \texttt{mimic\_LABEVENTS}
\end{itemize}

\textbf{OMOP\_Synthea}: Maps OMOP concepts to Synthea (synthetic patient generator) data. Example mapping:
\begin{itemize}
    \item \texttt{omop\_PERSON} $\rightarrow$ \texttt{synthea\_PATIENTS}
    \item \texttt{omop\_CONDITION\_OCCURRENCE} $\rightarrow$ \texttt{synthea\_CONDITIONS}
\end{itemize}

\section{Alternative Evaluation Methods}
\label{app:alt_eval}

In addition to the primary accuracy metrics (API-level and Field-level), we implement two alternative evaluation approaches that provide complementary perspectives on mapping quality.

\subsection{Embedding Similarity Evaluation}

We use sentence-transformer embeddings to compute semantic similarity between source and target field descriptions. This provides a fast, scalable metric independent of LLM judgments.

\textbf{Method:}
\begin{enumerate}
    \item Encode field descriptions using \texttt{paraphrase-MiniLM-L6-v2} sentence-transformer
    \item Compute cosine similarity between source and target embeddings
    \item Classify mapping as correct if similarity $> 0.3$ threshold
\end{enumerate}

\begin{equation}
\text{sim}(s, t) = \frac{\mathbf{e}_s \cdot \mathbf{e}_t}{\|\mathbf{e}_s\| \|\mathbf{e}_t\|}
\end{equation}

where $\mathbf{e}_s$ and $\mathbf{e}_t$ are the sentence-transformer embeddings of the source and target field descriptions respectively.

\textbf{Advantages:}
\begin{itemize}
    \item Fast and deterministic (no LLM calls required)
    \item Consistent across runs
    \item Scalable to large datasets
\end{itemize}

\textbf{Limitations:}
\begin{itemize}
    \item May miss semantic relationships beyond surface-level similarity
    \item Threshold selection (0.3) requires tuning per domain
\end{itemize}

\subsection{LLM-as-Judge Evaluation}

We employ an LLM as a judge to evaluate mappings on two semantic dimensions: equivalence and joinability. This captures domain-aware reasoning that embedding similarity may miss.

\textbf{Metrics:}
\begin{itemize}
    \item \textbf{Equivalence Score} (0--1): Do both fields serve the same functional purpose within their respective schemas?
    \item \textbf{Joinability Score} (0--1): Can one field's output be used to populate or enrich the other?
\end{itemize}

\textbf{Method:}
\begin{enumerate}
    \item For each predicted mapping $(s, t)$, retrieve field metadata from API lookup
    \item Construct a structured prompt with type names, field names, and descriptions
    \item Query WatsonX (Mistral-Large) to generate scores and explanation
    \item Classify mapping as correct if equivalence or joinability score $> 0.5$
\end{enumerate}

\textbf{Prompt Template:}
\begin{verbatim}
You are an API-mapping evaluator.
Your job is to judge whether two API schema
fields are functionally equivalent and/or
joinable based on their type names and
descriptions.

Endpoint A: {field_name_a} (type: {type_a})
Description: {description_a}

Endpoint B: {field_name_b} (type: {type_b})
Description: {description_b}

Return JSON with: equivalence_score,
joinability_score, explanation
\end{verbatim}

\textbf{Advantages:}
\begin{itemize}
    \item Captures semantic relationships beyond embedding similarity
    \item Provides interpretable explanations for each judgment
    \item Can assess joinability (data flow) in addition to equivalence
\end{itemize}

\textbf{Limitations:}
\begin{itemize}
    \item Requires LLM API calls (cost and latency)
    \item Non-deterministic (may vary across runs)
    \item Dependent on LLM capability and prompt engineering
\end{itemize}

\subsection{Evaluation Pipeline}

The complete evaluation pipeline combines all three approaches:

\begin{enumerate}
    \item \textbf{Primary Evaluation}: Calculate API-level and Field-level accuracy using exact match against ground truth (Appendix~\ref{app:accuracy})
    \item \textbf{Similarity Evaluation}: Compute embedding similarity scores for all predicted mappings
    \item \textbf{LLM-as-Judge}: Evaluate semantic equivalence and joinability for contested or low-confidence mappings
\end{enumerate}

Results are aggregated and reported as:
\begin{itemize}
    \item Accuracy tables by method/model/dataset (Tables~\ref{tab:api_accuracy}, \ref{tab:field_accuracy})
    \item Similarity score distributions
    \item LLM-as-judge positive rates (percentage scoring $> 0.5$)
\end{itemize}

\section{Reproducibility}
\label{app:reproducibility}

All experiments can be reproduced using the unified experiment runner:

\begin{verbatim}
# Run all experiments
./run.sh

# Run specific method/dataset/model
./run.sh --method map_reflect \
         --dataset ITOps --model granite

# Calculate accuracy metrics
./run.sh --accuracy

# Full pipeline: run + aggregate + accuracy
./run.sh --all
\end{verbatim}

\textbf{Evaluation Scripts:}
\begin{itemize}
    \item \texttt{eval/calc\_accuracy.py}: Primary accuracy calculation
    \item \texttt{eval/similarity\_evaluator.py}: Embedding similarity evaluation
    \item \texttt{eval/llm\_evaluator.py}: LLM-as-judge evaluation
    \item \texttt{eval/eval\_schema\_mapping.py}: Result aggregation and LaTeX generation
\end{itemize}

\textbf{Environment Requirements:}
\begin{itemize}
    \item Python 3.10+
    \item WatsonX API credentials (for LLM-based methods)
    \item ChromaDB for vector storage
    \item sentence-transformers for embeddings
\end{itemize}

\clearpage
\bibliographystyle{ACM-Reference-Format}
\bibliography{sample}

\end{document}
\endinput


